\setcounter{chapter}{0}
\chapter{Theoretical Study}

\graphicspath{{Chapitre1/figures/}}
%==============================================================================
\pagestyle{fancy}
\fancyhf{}
\fancyhead[R]{\bfseries\chaptername~\thechapter. }
\fancyfoot[R]{\thepage}
\renewcommand{\headrulewidth}{0.5pt}
\renewcommand{\footrulewidth}{0pt}
%\renewcommand{\chaptermark}[1]{\markright{\MakeUppercase{\chaptername~\thechapter. #1 }}{}}
%\renewcommand{\sectionmark}[1]{\markright{\thechapter.\thesection~ #1}}

\begin{spacing}{1.2}
%==============================================================================

\section*{Introduction}
Une étude théorique \cite{knuthwebsite} peut contenir l'une et/ou l'autre de ces deux parties :
\section{État de l'art} 
C'est une étude assez détaillée sur ce qui existe sur le marché ou dans la littérature (d'où 
le terme état de l'art), qui permet de répondre à la problématique. L'idée ici est de faire 
un comparatif entre les solutions existantes, mais surtout d'analyser le résultat de cette 
comparaison et de dire pourquoi ne sont-elles pas satisfaisantes pour répondre à votre 
problématique.
\begin{figure}[!ht]\centering
\includegraphics[scale=0.9]{art.jpg}
\caption{État de l'art}
\label{fig:fig1}
\end{figure}
\section{Étude de l'existant}
Elle est en général réalisée quand on va développer un module supplémentaire sur un 
logiciel existant, ou si on va modifier une application existante. L'étude de l'existant
consiste à expliquer ce qui existe déjà dans votre environnement de travail.

\section*{Conclusion}
La conclusion est en général sans numérotation, et n'apparaît pas dans la table des matières.


%==============================================================================
\end{spacing}
